\section{Verification and Reproducibility Details}
\label{app:verification}


\subsection{Candidate configurations}
\label{app:hyperparameters}

\noindent Table~\ref{tab:hyperparameters} records the compared model families
and the fixed settings that materially affect interpretation, including Logistic
Regression (LR), Random Forest (RF), XGBoost (XGB), Multilayer Perceptron (MLP),
and Principal Component Analysis (PCA). Figure~\ref{fig:feature_importances}
shows the relative importances of the handcrafted biophysical features in the
baseline XGBoost classifier.

\begin{figure}[!htbp]
\centering
\includegraphics[width=\linewidth]{../plots/fixed/feature_importances.png}
\caption{Handcrafted feature importances for the baseline XGBoost classifier.}
\label{fig:feature_importances}
\end{figure}


\begin{table}[!htbp]
\centering
\scriptsize
\setlength{\tabcolsep}{3pt}
\begin{tabular}{p{0.25\columnwidth}p{0.67\columnwidth}}
\toprule
Component & Candidates and fixed settings \\
\midrule
Handcrafted classifier & LR, RF, XGB; seven features; RF/XGB depth 3; grouped five-fold CV \\
Frozen/LoRA classifier & L1/L2 LR, RF, XGB, MLP; mean or mean+max+std pooling; RF 300 trees/depth 3; XGB 200 trees/depth 3; MLP after PCA-50 \\
Attention classifier & Additive token attention; masked pooling, dropout 0.4, weight decay $10^{-2}$, 120 epochs \\
Frozen/LoRA regressor & Ridge, RF, XGB; mean or mean+max+std; RF 300 trees/depth 4; XGB 200 trees/depth 3 \\
OAS adaptation & One rank-16 LoRA run; $\alpha=32$, dropout 0.05, learning rate $10^{-4}$, one epoch, effective batch 32 \\
Per-residue VAE & One run per representation; 1536--512--256--32 encoder, mirrored decoder, dropout 0.1, $\beta=0.005$, 25-epoch warm-up, 300 epochs \\
\bottomrule
\end{tabular}
\caption{Evaluated configurations. Because these candidates share reporting
folds, selected cross-validation estimates remain exploratory.}
\label{tab:hyperparameters}
\end{table}

\subsection{Classification detail}

\noindent Table~\ref{tab:classification_detail} reconciles the configuration
selection score with the pooled estimates used in the main text.

\begin{table}[!htbp]
\centering
\scriptsize
\setlength{\tabcolsep}{3pt}
\begin{tabular}{lrrr}
\toprule
Representation & AUROC & AUPRC & Fold AUROC \\
\midrule
Handcrafted XGB & 0.625 & 0.742 & $0.644\pm0.086$ \\
Frozen ESM3 XGB & 0.746 & 0.859 & $0.761\pm0.131$ \\
OAS LoRA ESM3 RF & 0.726 & 0.829 & $0.745\pm0.140$ \\
\bottomrule
\end{tabular}
\caption{Classification metrics reconciled from aligned OOF predictions.
Frozen XGB uses mean+max+std pooling; LoRA RF uses mean pooling. Configurations
maximize mean-fold AUROC.}
\label{tab:classification_detail}
\end{table}

\subsection{Regression detail}

\noindent The fixed, mean-fold-selected regression configurations give the
pooled estimates in Table~\ref{tab:regression_detail}.

\begin{table}[!htbp]
\centering
\small
\begin{tabular}{lrrr}
\toprule
Representation & Spearman & Pearson & $R^2$ \\
\midrule
Frozen ESM3 & 0.366 & 0.357 & 0.114 \\
OAS LoRA ESM3 & 0.319 & 0.310 & 0.036 \\
Length-only Ridge & 0.004 & $-0.059$ & $-0.046$ \\
\bottomrule
\end{tabular}
\caption{Pooled OOF breadth-regression metrics. Frozen and LoRA configurations
maximize mean-fold Spearman.}
\label{tab:regression_detail}
\end{table}

\subsection{Lineage-group leakage audit}
\label{app:grouping_audit}

\noindent The fallback grouping (Section~\ref{sec:data}) does not cross-merge
tiers, so an antibody can resemble a member of a different group. We surface
every such case with \texttt{find\_cross\_group\_near\_duplicates}: across all
$\binom{300}{2}$ antibody pairs, nine span a group boundary while exceeding the
$80\%$ Levenshtein-identity threshold (Table~\ref{tab:grouping_audit}). Several
involve inferred ancestral sequences (unmutated common ancestors and
intermediates, e.g.\ \texttt{VRC34-UCA}, \texttt{DH270 UCA/IA4}) or
germline-reverted variants that share a $V_H$ germline across lineages or donors.
They are deliberately \emph{not} auto-merged; the count quantifies the residual
leakage left after grouping rather than eliminating it silently.

\begin{table}[!htbp]
\centering
\scriptsize
\setlength{\tabcolsep}{3pt}
\begin{tabular}{llp{0.20\columnwidth}p{0.20\columnwidth}r}
\toprule
Antibody A & Antibody B & Group A & Group B & Id. \\
\midrule
1361      & D5\_AR        & patient:donor\_Gorny2012 & cluster:10       & 0.811 \\
1361      & PGZL1\_gVmDmJ & patient:donor\_Gorny2012 & lineage:PGZL1    & 0.839 \\
1393A     & D5\_AR        & patient:donor\_Gorny2012 & cluster:10       & 0.803 \\
1393A     & PGZL1\_gVmDmJ & patient:donor\_Gorny2012 & lineage:PGZL1    & 0.855 \\
D5\_AR    & PGZL1\_gVmDmJ & cluster:10               & lineage:PGZL1    & 0.815 \\
D5\_FI    & PGZL1\_gVmDmJ & cluster:10               & lineage:PGZL1    & 0.806 \\
DH270 IA4 & VRC34-UCA     & lineage:DH270            & lineage:VRC34    & 0.843 \\
DH270 UCA & VRC34-UCA     & lineage:DH270            & lineage:VRC34    & 0.874 \\
PGDM1400  & PGDM1401      & lineage:PGDM1400         & patient:Donor 84 & 0.816 \\
\bottomrule
\end{tabular}
\caption{Residual cross-group near-duplicates: antibody pairs with
$\geq 80\%$ Levenshtein identity assigned to different lineage groups. Group
labels carry their fallback tier (\texttt{lineage}, \texttt{patient},
\texttt{cluster}). These pairs are intentionally not merged, so the table
measures the leakage remaining after the hierarchical grouping.}
\label{tab:grouping_audit}
\end{table}

\subsection{Maturation correlation sensitivity}

\noindent Small lineage sizes make the Fisher-$z$ sensitivity intervals in
Table~\ref{tab:tier2} wide, especially for DH270 and DH511.

\begin{table}[!htbp]
\centering
\scriptsize
\setlength{\tabcolsep}{3pt}
\begin{tabular}{lrrr}
\toprule
Lineage & $n$ & Frozen $r$ [interval] & LoRA $r$ [interval] \\
\midrule
VRC34 & 14 & 0.651 [0.184, 0.878] & 0.923 [0.768, 0.976] \\
VRC01 & 15 & 0.759 [0.404, 0.915] & 0.960 [0.881, 0.987] \\
DH270 & 6  & 0.771 [$-0.108$, 0.973] & 0.995 [0.950, 0.999] \\
DH511 & 8  & 0.084 [$-0.660$, 0.745] & 0.409 [$-0.415$, 0.865] \\
\bottomrule
\end{tabular}
\caption{Descriptive Pearson correlations and Fisher-$z$ sensitivity intervals (maturation correlation analysis). These are not independent estimates of a causal adaptation effect.}
\label{tab:tier2}
\end{table}

\subsection{Interpretation rules}

The paper uses ``interval'' rather than ``confidence interval'' where the
resampling assumptions are especially limited. It avoids significance claims for
the VAE evaluations, does not treat non-overlap or interval sign as proof of a
biological mechanism, and distinguishes direct observations from causal
interpretation. Project history retains random-split and earlier name-grouped
scores, but the main result narrative excludes them because their
current prediction artifacts are not part of the verified canonical bundle.

\subsection{Per-configuration discriminative results}

Tables~\ref{tab:all_classification} and~\ref{tab:all_regression} report every
fixed-pooling configuration compared for the frozen and LoRA representations.
The entries are not independent experiments: all candidates reuse the same five
reporting folds. Selection uses the mean-fold column, even when a different
configuration has a larger pooled estimate. The complete classification
performance across all heads and pooling schemes appears in
Figure~\ref{fig:all_classification_auroc}, while
Figure~\ref{fig:all_regression_spearman} details the regression Spearman
correlations.

\begin{figure}[!htbp]
\centering
\includegraphics[width=\linewidth]{../plots/lora/classification_auroc_by_head.png}
\caption{Mean grouped-fold AUROC (with standard deviation) for all evaluated classifier heads and pooling configurations under frozen and LoRA representations.}
\label{fig:all_classification_auroc}
\end{figure}

\begin{figure}[!htbp]
\centering
\includegraphics[width=\linewidth]{../plots/lora/regression_spearman_by_config.png}
\caption{Regression Spearman correlation across different model configurations, comparing frozen and LoRA representations.}
\label{fig:all_regression_spearman}
\end{figure}

\begin{table}[!htbp]
\centering
\scriptsize
\setlength{\tabcolsep}{3pt}
\begin{tabular}{lrrrr}
\toprule
& \multicolumn{2}{c}{Frozen} & \multicolumn{2}{c}{LoRA} \\
Configuration & AUROC & AUPRC & AUROC & AUPRC \\
\midrule
XGB, mean+max+std & 0.761 & 0.878 & 0.668 & 0.794 \\
RF, mean+max+std & 0.759 & 0.874 & 0.740 & 0.855 \\
RF, mean & 0.759 & 0.875 & \textbf{0.745} & 0.858 \\
XGB, mean & 0.744 & 0.864 & 0.692 & 0.808 \\
L1 LR, mean+max+std & 0.732 & 0.858 & 0.667 & 0.805 \\
L2 LR, mean+max+std & 0.702 & 0.836 & 0.670 & 0.797 \\
L1 LR, mean & 0.697 & 0.827 & 0.741 & 0.834 \\
MLP, mean & 0.696 & 0.836 & 0.700 & 0.803 \\
L2 LR, mean & 0.677 & 0.822 & 0.723 & 0.833 \\
MLP, mean+max+std & 0.668 & 0.804 & 0.679 & 0.783 \\
\bottomrule
\end{tabular}
\caption{Mean grouped-fold classification results for all frozen and LoRA
fixed-pooling candidates. Bold marks the selected LoRA configuration; the
selected frozen configuration is the first row.}
\label{tab:all_classification}
\end{table}

\begin{table}[!htbp]
\centering
\scriptsize
\setlength{\tabcolsep}{3pt}
\begin{tabular}{lrrrr}
\toprule
& \multicolumn{2}{c}{Frozen Spearman} &
\multicolumn{2}{c}{LoRA Spearman} \\
Configuration & Fold & Pooled & Fold & Pooled \\
\midrule
RF, mean+max+std & \textbf{0.380} & 0.366 & 0.304 & 0.276 \\
RF, mean & 0.342 & 0.344 & \textbf{0.357} & 0.319 \\
XGB, mean+max+std & 0.341 & 0.354 & 0.292 & 0.272 \\
XGB, mean & 0.329 & 0.340 & 0.353 & 0.333 \\
Ridge, mean+max+std & 0.200 & 0.233 & 0.226 & 0.235 \\
Ridge, mean & 0.149 & 0.171 & 0.208 & 0.236 \\
Length-only Ridge & 0.155 & 0.004 & 0.155 & 0.004 \\
\bottomrule
\end{tabular}
\caption{Mean-fold and pooled-OOF regression Spearman correlations for all
candidates. Bold marks the configurations selected by mean-fold Spearman.}
\label{tab:all_regression}
\end{table}

\subsection{Paired resampling results}

The verification script aligns predictions by antibody ID and resamples complete
\texttt{group\_id} units. Table~\ref{tab:paired_deltas} gives the resulting
paired differences; 10,000 valid resamples produce every interval.

\begin{table}[!htbp]
\centering
\scriptsize
\setlength{\tabcolsep}{4pt}
\begin{tabular}{llrrr}
\toprule
Comparison & Metric & $\Delta$ & Low & High \\
\midrule
Frozen--handcrafted & AUROC & 0.121 & 0.012 & 0.228 \\
Frozen--handcrafted & AUPRC & 0.117 & 0.028 & 0.207 \\
LoRA--frozen & AUROC & $-0.020$ & $-0.096$ & 0.059 \\
LoRA--frozen & AUPRC & $-0.031$ & $-0.075$ & 0.024 \\
LoRA--frozen & Spearman & $-0.047$ & $-0.169$ & 0.064 \\
LoRA--frozen & Pearson & $-0.047$ & $-0.173$ & 0.071 \\
LoRA--frozen & $R^2$ & $-0.078$ & $-0.228$ & 0.053 \\
\bottomrule
\end{tabular}
\caption{Paired group-bootstrap differences for the selected configurations.
Intervals quantify sensitivity to the observed lineage groups and do not
correct for configuration selection.}
\label{tab:paired_deltas}
\end{table}

\subsection{VAE reconstruction and geometry summary}

The two VAEs use identical split membership and architecture, but we train them
in separate single runs. The frozen and LoRA best epochs are 299 and 296,
respectively; their best validation total/reconstruction/KL losses are
$0.2599/0.1859/14.80$ and $0.2834/0.2045/15.79$. Each run uses 27,490 training,
5,738 validation, and 5,901 held-out residue embeddings. The detailed VAE
geometry checks and verification metrics appear in Table~\ref{tab:vae_verification}.

\begin{table}[!htbp]
\centering
\scriptsize
\setlength{\tabcolsep}{4pt}
\begin{tabular}{lrr}
\toprule
Verified VAE quantity & Frozen & LoRA \\
\midrule
Mean Pearson correlation (maturation analysis) & 0.566 & 0.822 \\
Mean Spearman correlation (maturation analysis) & 0.384 & 0.692 \\
Mean VAE--line MAE (geometry check) & 0.041 & 0.076 \\
Standard deviation (geometry check) & 0.274 & 0.237 \\
Median (geometry check) & 0.004 & 0.018 \\
Positive fraction (geometry check) & 0.573 & 0.693 \\
Lineage bootstrap low (geometry check) & $-0.016$ & 0.023 \\
Lineage bootstrap high (geometry check) & 0.109 & 0.142 \\
\bottomrule
\end{tabular}
\caption{Descriptive VAE geometry summaries. Positive VAE--line
differences (geometry check) favor the reconstructed straight-line control.}
\label{tab:vae_verification}
\end{table}


\subsection{Reproducibility limits}

The CPU verification script reconstructs claims from saved OOF predictions and
evaluation artifacts; it does not retrain ESM3, the LoRA adapter, or either
VAE\@. Therefore, it verifies internal consistency and reported calculations,
not bitwise reproducibility of the GPU training runs. Similarly, the group
bootstrap resamples the observed 113 metadata-derived groups. It does not
represent a new external cohort, and its intervals remain conditional on the
selected model configurations.

Confirmatory evidence would require a locked external lineage cohort or nested
grouped model selection, repeated adapter and VAE seeds, and raw-embedding
directionality controls for every lineage. Those missing experiments are
limitations of the current evidence, not tasks silently replaced by the claim
verification.
